%*************************************************
% In this file the abstract is typeset.
% Make changes accordingly.
%*************************************************


\addcontentsline{toc}{section}{چکیده}
\newgeometry{left=2.5cm,right=3cm,top=3cm,bottom=2.5cm,includehead=false,headsep=1cm,footnotesep=.5cm}
\setcounter{page}{1}
\pagenumbering{arabic}						% شماره صفحات با عدد
\thispagestyle{empty}


\subsection*{چکیده}
\begin{small}
\baselineskip=0.7cm

یادگیری ماشین\LTRfootnote{Machine Learning} به‌عنوان یکی از شاخه‌های برجسته هوش مصنوعی، در دهه‌های اخیر با رشد چشمگیری در صنایع مختلف از جمله کشاورزی، سلامت، و تولید مواجه شده است. این فناوری با بهره‌گیری از الگوریتم‌های پیشرفته، امکان تحلیل داده‌های پیچیده را فراهم کرده و تصمیم‌گیری‌های مبتنی بر داده را بهبود بخشیده است. در صنعت دامداری، داده و استفاده از روابط داده‌ها نقش کلیدی در بهینه‌سازی تولید و مدیریت گله‌ها ایفا می‌کند و با افزایش حجم داده‌های ثبت‌شده، روند استفاده از این روش‌ها از مدل‌های آماری سنتی به سمت روش‌های هوشمند مبتنی بر یادگیری ماشین سوق یافته است. 

دوره انتقال\LTRfootnote{Transition period}، شامل سه هفته پیش و سه هفته پس از زایش گاوهای شیری، دوره‌ای حساس است که تغییرات فیزیولوژیکی و متابولیکی گسترده‌ای را به همراه دارد. این تغییرات می‌توانند بر سلامت گاو و موفقیت دوره شیردهی تأثیر بگذارند. در این میان، شاخص دوره انتقال\LTRfootnote{Transition Cow Index (TCI)} به‌عنوان ابزاری برای ارزیابی غیرمستقیم مدیریت گله در دوره انتقال معرفی شده است. این شاخص از تفاوت مقدارهای پیش‌بینی‌شده و مقدار واقعی برای اولین رکورد شیردهی به دست می‌آید و انحراف به‌دست‌آمده به‌عنوان معیاری برای کیفیت مدیریت در این دوره بحرانی به کار می‌رود. با این حال، روش‌های مستقیم مانند پروفایل خونی گاو نیز وجود دارند؛ اما به دلیل خطاهایی که ممکن است داشته باشند و هزینه‌های بالای آن‌ها، استفاده از این روش‌ها محدود شده است. هدف این پژوهش، توسعه و مقایسه مدل‌های پیش‌بینی برای تخمین اولین رکورد شیردهی گاوها با استفاده از داده‌های دوره‌های شیردهی پیشین\LTRfootnote{Previous Lactation} و داده‌های روزهای تست\LTRfootnote{Test Days} است. این کار به منظور ارائه ابزاری نوین برای ارزیابی مدیریت دوره انتقال انجام شده است.

برای این منظور، مجموعه‌ای جامع از داده‌های واقعی مربوط به دامداری‌های کل کشور جمع‌آوری و سازماندهی شده است. این داده‌ها از طریق همکاری با تعاونی دامداران وحدت استان اصفهان، به‌عنوان یک مرکز داده ملی، گردآوری شده‌اند. این مجموعه داده شامل اطلاعات 345,676 رأس گاو و 99 گله است که از سال 1390 تا 1401 جمع‌آوری شده‌اند.

مدل‌های پیش‌بینی شامل رگرسیون ریج اولیه\LTRfootnote{Ridge Regression}، رگرسیون ریج بهینه‌شده، جنگل تصادفی اولیه\LTRfootnote{Random Forest}، جنگل تصادفی بهینه‌شده، جنگل تصادفی با اعتبارسنجی متقاطع نهایی،  جنگل تصادفی اولیه با استفاده از گروه‌بندی، جنگل تصادفی بهینه‌شده با استفاده از گروه‌بندی، جنگل تصادفی با استفاده از گروه‌بندی و اعتبارسنجی متقاطع نهایی، رگرسیون بردار پشتیبانی\LTRfootnote{Support Vector Regression(SVR)} با هسته\LTRfootnote{Kernel} خطی\LTRfootnote{Linear}، رگرسیون بردار پشتیبانی با هسته چندجمله‌ای\LTRfootnote{Polynomial}، رگرسیون بردار پشتیبانی با هسته شعاعی\LTRfootnote{RBF}، رگرسیون بردار پشتیبانی با هسته شعاعی و با مقیاس پذیری مقاوم داده ها، افزایش گرادیان افراطی\LTRfootnote{XGBoost}، افزایش گرادیان افراطی بهینه‌شده، تقویت دسته‌ای اولیه\LTRfootnote{CatBoost}، تقویت دسته‌ای بهینه‌شده، تقویت دسته‌ای با مقیاس پذیری استاندارد داده ها، شبکه عصبی چندلایه\LTRfootnote{MultiLayer Perceptron(MLP)}، و شبکه عصبی عمیق\LTRfootnote{Deep Neural Network(DNN)} هستند که برای پیش‌بینی اولین رکورد شیردهی به کار گرفته شده‌اند. عملکرد این مدل‌ها با استفاده از معیارهای ضریب تعیین\LTRfootnote{R-squared (R²)}، میانگین قدرمطلق خطا\LTRfootnote{Mean Absolute Error (MAE)}، ریشه میانگین مربع خطا\LTRfootnote{Root Mean Squared Error (RMSE)}، درصد خطای میانگین\LTRfootnote{Mean Percentage Error (MPE)}، میانگین قدرمطلق درصد خطای متقارن\LTRfootnote{Symmetric Mean Absolute Percentage Error (sMAPE)}، و انحراف استاندارد نسبی\LTRfootnote{Standard Deviation Ratio (SDR)} ارزیابی شده‌اند. 

نتایج نشان داد که مدل تقویت دسته‌ای اولیه، بهترین عملکرد را با ضریب تعیین \num{0.4015} و میانگین قدرمطلق خطا \num{6.4382} داشته است.

عملکرد و دقت مدل تقویت دسته‌ای اولیه با ضریب تعیین و میانگین قدرمطلق خطا، نسبت به مدل فعلی مورد استفاده در تعاونی وحدت استان اصفهان که ضریب تعیین حدود \num{0.30} دارد، تقریباً 10 درصد بهبود در دقت پیش‌بینی را نشان می‌دهد. همچنین، این مدل در مقایسه با سایر مدل‌های ارزیابی‌شده، بهترین عملکرد را داشته است. برای مثال، نسبت به مدل شبکه عصبی چندلایه با ضریب تعیین \num{0.37}، حدود \num{3.14} درصد بهبود در ضریب تعیین دارد. در مقایسه با مدل جنگل تصادفی بهینه‌شده با \num{0.3575}، بهبود حدود \num{4.14} درصد، و نسبت به رگرسیون بردار پشتیبانی با هسته شعاعی با \num{0.3643}، بهبود حدود \num{3.72} درصد مشاهده می‌شود. این مدل همچنین نسبت به مدل افزایش گرادیان افراطی بهینه‌شده با \num{0.373}، حدود \num{2.85} درصد عملکرد بهتری ارائه می‌دهد. این نتایج نشان‌دهنده برتری قابل‌توجه مدل تقویت دسته‌ای اولیه نسبت به سایر مدل‌های بررسی‌شده و مدل فعلی تعاونی وحدت استان اصفهان است.

برای تحقیقات آینده، پیشنهاد می‌شود تمرکز بر استفاده و بهینه‌سازی الگوریتم‌های پیشرفته‌تر شبکه‌های عصبی مصنوعی، مانند مدل‌های عمیق‌تر، باشد تا دقت پیش‌بینی و انعطاف‌پذیری مدل‌ها در شرایط مختلف بهبود یابد.
\vspace{0.5cm}

\noindent\textbf{کلمات کلیدی:} دوره انتقال، مدیریت دامداری، یادگیری ماشین، پیش‌بینی شیر، اولین رکورد شیردهی\LTRfootnote{First Test Day Milk Yield}

\end{small}